\documentclass[a4paper,12pt]{article}
\usepackage{ragged2e}
\usepackage{graphicx}
\usepackage[
    colorlinks
]{hyperref}

\hypersetup{
    pdfinfo = {
        Title = Public Opinion and Mass Media,
        Author = {Saad Nawaz Ghauri},
        Subject = Mass Communication assignment,
        Keywords = {mass media, public opinion, propaganda},
        Producer = LuaLaTeX
    },
    bookmarks = true,
    bookmarksopen = true,
    pdfmenubar = true,
    pdfnewwindow = true,
    pdftoolbar = true,
    verbose = true,
}

\author{Saad Nawaz Ghauri}
\title{Public Opinion and Mass Media}

\begin{document}
\maketitle

\begin{Center}
\includegraphics{nuces_logo.png}
\end{Center}

\begin{Center}
    Section: \emph{BCS-4A} \\
    Role Number: \emph{23L-1007} \\
    Course Code: \emph{SS-2040} \\
    Course Title: \emph{Mass Communication} \\
\end{Center}

\begin{Center}
    School of Computing \\
    National University of Computer and Emerging Sciences, \\
    Lahore Campus \\
\end{Center}

\newpage

\tableofcontents
\newpage


\section{Introduction}
This article will explore the complex relationship between mass media and public opinion, how they influence each other.

\section{Definition}
Before attempting to understand or measure the way public opinion and mass media interact, it would be useful to define what these terms are.
\subsection{Public Opinion}
Public opinion is the aggregate of the views, opinions, attitudes, and beliefs held by the majority of the public,
when it comes to a topic, or at least by a significant part of soceity.
\subsection{Mass Media}

\section{Main Theme}
\subsection{Positive Effects}
\subsection{Negative Effects}
\subsection{Examples}

\section{Conclusion}


\end{document}
